\section{局部稳定性、吸引盆与分离界}

\begin{definition}[稳定状态]
如果从某个状态足够近处出发的轨迹，倾向于停留在其附近，或在小扰动后重新返回其附近，
则该状态是\textbf{局部稳定}的。
\end{definition}

\begin{definition}[盆地与分离界]
稳定状态的\textbf{吸引盆}，是指所有会流向该状态的初始条件集合。
\textbf{分离界}（separatrix）则是状态空间中的边界，用来划分那些会走向不同长期结果的轨迹。
\end{definition}

\begin{discussion}
这些定义并非装饰性的。
它们为“行为失败”提供了更锐利的解释。
一个人并不总是因为“诱惑太大”而失败。
有时只是因为系统起始时已经深处错误的吸引盆之内，
或者距离分离界太远，以至于一个小干预不足以改变方向。
这正是时机与准备如此重要的原因：
它们决定了系统是在仍靠近边界时被推动，还是在已经深陷自我强化区域之后才被处理。
\end{discussion}

\begin{example}[休息后重新进入学习，作为一个分离界问题]
设一名学生吃完午饭回到书桌前。
如果课本已经翻开、下一题已经标出、手机不在场，那么午休后的状态可能正靠近“学习盆地”和“分心盆地”之间的边界。
此时一个很小的线索就足以重新导向轨迹。
如果桌面凌乱且手机已经激活，同一个学生就可能一开始便深处于分心盆地内部。
两者差异不在道德价值，而在初始条件的几何位置。
\end{example}